% -----------------------------------------------------------------------------------------------------
\section{Glauber 有效理论、旁观者相消与因子化破坏}\label{sec:front-glauber}

单举 Drell--Yan 中 Glauber 区是领头幂次的，但在对割求和与围道形变之后相消\cite{Bodwin1985,CollinsSoperSterman1985b,CollinsSoperSterman1988}（\cref{thm:dy-factorization}，[A]，瓶颈见附录~O 的假设登记）。Collins 在“遗留问题”一节\cite[第~14.4.8~节，第~561--562~页]{Collins2011}指出，按快度分级的交换“very quickly”把人带到“trying to develop a good and deductive QCD version of Regge theory”的境地，并说“the current versions of genuine proofs are evidently formidable”。本节记录 2010 年以后的进展；它们都不改变\cref{thm:dy-factorization}的等级。

\begin{proposition}[Glauber 有效理论：Regge 化、BFKL 与零箱减除 \tagA]\label{prop:front-glauber-bfkl}
在加入\cref{prop:scet-glauber-operators}的 Glauber 势算符的 SCET 中（领头幂次拉氏量\cite[式~(5.21)]{RothsteinStewart2016}已在该处列出，这里沿用第十五章的 SCET 约定）：(i) 算符基由树图与单圈匹配确定，且“there are no loop corrections to this form”\cite[第~5.3~节]{RothsteinStewart2016}；(ii) 单圈前向振幅的快度重整化给出胶子 Regge 化，前向散射中领头软函数与共线函数的快度重整化群方程就是 BFKL 方程（\cref{thm:smallx-bfkl}）；(iii) Glauber 零箱减除使硬散射中的软、共线圈图对 Wilson 线方向不敏感，反之部分 Glauber 交换可以通过选择 Wilson 线方向吸收进 Wilson 线——这是 CSS“形变进 Wilson 线”的有效理论版本。软算符 $\mathcal O_s^{BC}$ 的单胶子 Feynman 规则是 Lipatov 顶点，它也支配两个以上软胶子的发射。先驱工作指出，为使旁观者相互作用自洽，SCET 必须加入 Glauber 模式\cite{Bauer:2010cc}（据该文摘要）。（本命题只涉及部分子散射。）
% VF-FRO-6
% errtype: E-S
\uses{prop:scet-glauber-operators, thm:smallx-bfkl, def:scet-zero-bin}
\end{proposition}

\begin{proofsketch}
从 QCD 到 SCET$_{\rm I,II}$ 的树图与单圈匹配确定算符；对单圈前向振幅作快度正规化与重整化，读出 Regge 轨迹与 BFKL 核；零箱减除的论证在单圈与两圈显式进行。完整论证见\cite[第~5、9--11~节]{RothsteinStewart2016}。
\end{proofsketch}

\paragraph{地位与缺口。}[A]：与已知结果一致，(i) 是论证而非证明。Rothstein 与 Stewart 自己说，这一构造“taken alone it is far from a proof of factorization, even in the abelian case”；他们对 $e^+e^-$ 中的 $n$--$\bar n$ 产生流在单圈与两圈处理了 Glauber 贡献（在 SCET$_{\rm II}$ 中显式出现并等于软零箱，或在 SCET$_{\rm I}$ 中由超软图承担），并期望任意阶同样如此，但没有证明；BPS 场重定义并不消去 Glauber 耦合\cite[第~5.1~节]{RothsteinStewart2016}。“Rothstein--Stewart（2016）证明了 Drell--Yan 因子化”是过度表述。旁观者—旁观者相位 $\ee^{\ii\phi(b_\perp)}$ 当且仅当 $\Delta p_\perp$ 积分不受限制时相消（\cite[式~(9.1)、(9.18)--(9.19)、(11.12)--(11.13)]{RothsteinStewart2016}，\cite[式~(14)]{AybatSterman2009}），见第十二章（\cref{prop:dy-abelian-model}）。

\begin{proposition}[纯 Glauber 梯形图在单标度红外安全观测量中的相消 \tagA]\label{prop:front-syz}
对纯 Glauber 梯形图（旁观者—旁观者交换，旁观者微扰），振幅层次一切破坏因子化的效应在截面层次对\emph{任何}红外安全的单标度观测量（如 $E_T$、束流推力）相消：$\int\dd z\,\dd\bar z\sum_{\rm cuts}(\text{梯形图})=0$。推论：领头幂次的破坏需要从 Glauber 梯级上辐射软胶子（Lipatov 顶点），这把因子化破坏与底层事件联系起来\cite{Schwartz:2018obd}。（本命题只涉及微扰旁观者。）
% errtype: E-S
\uses{prop:front-glauber-bfkl, prop:scet-glauber-operators}
\end{proposition}

\begin{proofsketch}
用共形变量 $(z,\bar z)$ 描述旁观者的横向运动学；由 $2\to2$ 振幅的标度不变性，对每个单标度观测量，相空间积分对 $z$ 是单举的。NLO 显式核对：两圈的 $O(\epsilon^0)$ 项与单圈的 $\epsilon\times1/\epsilon$ 项相消。全阶：梯形图的红外发散指数化为相位（其奇异部分见原文式~(53)），余项可积。完整论证见\cite{Schwartz:2018obd}。
\end{proofsketch}

\paragraph{地位与缺口。}[A]，限于所述的一类图与观测量；原文的适用范围：“Our argument does not apply to doubly-differential observables, like those studied in [Zeng]”。分歧（并列给出，不作裁决）：Rothstein--Stewart 第~11.1~节称，当测量约束 $\Delta p_\perp\sim Q_*$ 且旁观者微扰时，旁观者—旁观者 Glauber 交换对束流推力、横向推力与 $E_T$“spoil the factorization \dots\ starting at $O(\alpha_s^4)$”；Schwartz、Yan 与 Zhu 对纯梯形图与单标度观测量的论证与之相反，并称其“disproves previous claims”——这一结论未被反驳，并为后来的工作采用；前者的一般机制只对多重微分的测量仍未被否定。

\begin{conjecture}[对多重部分子相互作用敏感的观测量：标准因子化的破坏 \tagC]\label{conj:front-mpi-violation}
对强子—强子碰撞中对多重部分子相互作用敏感的单标度观测量（$E_T$、束流推力 $B^+_{a,b}$、横向推力），标准的硬—软—共线因子化在领头幂次失效，机制是 Glauber 梯级上的软辐射，即多重部分子相互作用（底层事件）的贡献。
\uses{prop:front-syz}
\end{conjecture}

\paragraph{地位与缺口。}[C]。证据：Gaunt\cite{Gaunt2014}（据其摘要与\cite{Schwartz:2018obd}的转述）发现 CSS 的相消在两个旁观者之间交换两个 Glauber 胶子的层次上失效，并将其与额外的软散射联系起来——这是论证而非证明，其纯 Glauber 部分已被\cref{prop:front-syz}否定；SYZ 的推论指出了剩下的机制。对 \emph{广义}因子化（两个共线扇区解耦），带喷注否决或束流推力的 Drell--Yan 有一个模型反例：解耦意味着与强子纵向自旋无关，而旁观者—旁观者 Glauber 交换给出非零的领头幂次自旋不对称（含标量夸克的 QCD；\cite{Zeng:2015iba}，据该文摘要）——这是模型中的 [V]，核心陈述为\cref{prop:break-jet-veto}。不提\cite{Schwartz:2018obd}而说“Gaunt（2014）证明了强子 $E_T$ 的因子化破坏”是过度表述。

\begin{proposition}[重末态与旁观者之间的软末态相互作用 \tagA]\label{prop:front-heavy-fsi}
高 $p_T$ 重粒子与旁观者之间的软末态相互作用，对反冲伙伴不受约束的单粒子单举截面（如 $t\bar t$ 中单个顶夸克的单举分布）是幂次压低的，而对没有硬反冲的多粒子单举截面（包括粒子对的 $q_T$ 分布）给出领头幂次效应；非微扰效应压低 $\Lambda/p'_T$\cite{Mitov:2012gt}（据该文摘要）。相关地，$t\bar t$ 的小 $q_T$ 重求和需要带方位角关联的非平凡软因子\cite{Catani:2014qha}（据摘要）。（本命题只涉及部分子散射。）
% errtype: E-S
\uses{prop:method-fsi}
\end{proposition}

\begin{proofsketch}
在重末态与旁观者之间的软交换中，对未观测的反冲伙伴求和使末态相互作用由幺正性相消（与\cref{prop:method-fsi}同一机制）；观测多个粒子时这一求和不再完整，留下领头幂次效应。完整论证见\cite{Mitov:2012gt}。
\end{proofsketch}

\paragraph{地位与缺口。}[A]（只读到摘要层次）。

\begin{remark}[Glauber 现状与本书分级的后果]\label{rem:front-glauber-status}
单举 Drell--Yan 中的 Glauber 相消：[A]，不变（\cref{thm:dy-factorization}）。Rothstein--Stewart 拉氏量：[A]（\cref{prop:scet-glauber-operators}）。旁观者相位：对所述的程函计算为 [R]（第十二章）。纯梯形图的相消：[A]（\cref{prop:front-syz}）。对多重部分子相互作用敏感的观测量：[C]（\cref{conj:front-mpi-violation}）。带自旋不对称的广义因子化：模型中 [V]。严格的类空共线因子化：单圈 [V]\cite{CataniDeFlorianRodrigo2012}（\cref{prop:break-spacelike}）——他们没有证明任何相消，并说该效应“may invalidate the factorization theorem of mass singularities starting at the N$^3$LO”。非单举 $pp$ 截面的共线因子化：纯 QCD 中从三圈起（论证中有一步未证明；两圈效应在纯 QCD 中为零）、超出纯 QCD 从两圈起 [V]\cite{ForshawSeymourSiodmok2012}（\cref{prop:break-fss}），并已在带 Glauber 算符的 SCET 中重现\cite{Schwartz:2017nmr}。全色两圈软流中，软胶子与入射部分子共线时出现“潜在危险”的项\cite{Dixon:2019lnw}，但按其勘误（JHEP 06 (2024) 143），这些项在截面中对色求和后相消，不能作为截面层次破坏的证据。类时劈裂振幅已算到三圈，类空的只对三部分子振幅给出\cite{Guan:2024hlf}。误解 M21：严格共线因子化的破坏是关于类空多部分子极限的振幅层次陈述，它本身并不意味着 PDF 因子化失效。
\end{remark}
